\chapter{Optimization}

\section{Search on Answer}
\setcounter{section}{1}
\setcounter{subsection}{0}
\subsection{Binary Search}
Binary search can be generally used to find the input value corresponding to any output value of a monotonic (strictly non-increasing or strictly non-decreasing) function in $O(\log n)$ time with respect to the domain size. This is a special case of finding the exact point at which any given monotonic Boolean function changes from true to false or vice versa. Unlike searching through an array, discrete binary search is not restricted by available memory, making it useful for handling infinitely large search spaces such as real number intervals.

\begin{compactitem}
  \item \inlinecode{binary\_search\_first\_true()} takes two integers \inlinecode{lo} and \inlinecode{hi} as boundaries for the search space \inlinecode{[lo, hi)} (i.e. including \inlinecode{lo}, but excluding \inlinecode{hi}) and returns the smallest integer \inlinecode{k} in \inlinecode{[lo, hi)} for which the predicate \inlinecode{pred(k)} tests true. If \inlinecode{pred(k)} tests false for every \inlinecode{k} in \inlinecode{[lo, hi)}, then \inlinecode{hi} is returned. This function must be used on a range in which there exists a constant \inlinecode{k} such that \inlinecode{pred(x)} tests false for every \inlinecode{x} in \inlinecode{[lo, k)} and true for every \inlinecode{x} in \inlinecode{[k, hi)}.
  \item \inlinecode{binary\_search\_last\_true()} takes two integers \inlinecode{lo} and \inlinecode{hi} as boundaries for the search space \inlinecode{[lo, hi)} (i.e. including \inlinecode{lo}, but excluding \inlinecode{hi}) and returns the largest integer \inlinecode{k} in \inlinecode{[lo, hi)} for which the predicate \inlinecode{pred(k)} tests true. If \inlinecode{pred(k)} tests false for every \inlinecode{k} in \inlinecode{[lo, hi)}, then \inlinecode{hi} is returned. This function must be used on a range in which there exists a constant \inlinecode{k} such that \inlinecode{pred(x)} tests true for every \inlinecode{x} in \inlinecode{[lo, k]} and false for every \inlinecode{x} in \inlinecode{(k, hi)}.
  \item \inlinecode{fbinary\_search()} is the equivalent of \inlinecode{binary\_search\_first\_true()} on floating point predicates. Since any interval of real numbers is dense, the exact target cannot be found due to floating point error. Instead, a value that is very close to the border between false and true is returned. The precision of the answer depends on the number of repetitions the function performs. Since each repetition bisects the search space, the absolute error of the answer is $1/(2^n)$ times the distance between \inlinecode{lo} and \inlinecode{hi} after n repetitions. Although it is possible to control the error by looping while \inlinecode{hi - lo} is greater than an arbitrary epsilon, it is simpler to let the loop run for a desired number of iterations until floating point arithmetic break down. 100 iterations is usually sufficient, since the search space will be reduced to $2^{-100}$ (roughly $10^{-30}$) times its original size. This implementation can be modified to find the "last true" point in the range by simply interchanging the assignments of \inlinecode{lo} and \inlinecode{hi} in the if-else statements.
\end{compactitem}

\textbf{Time Complexity}:
\begin{compactitem}
  \item $O(\log n)$ calls will be made to \inlinecode{pred()} in \inlinecode{binary\_search\_first\_true()} and \inlinecode{binary\_search\_last\_true()}, where $n$ is the distance between \inlinecode{lo} and \inlinecode{hi}.
  \item $O(n)$ calls will be made to \inlinecode{pred()} in \inlinecode{fbinary\_search()}, where $n$ is the number of iterations performed.
\end{compactitem}

\textbf{Space Complexity}: $O(1)$ auxiliary for all operations.

\begin{lstlisting}[language=C++]
template<class Int, class IntPredicate>
Int binary_search_first_true(Int lo, Int hi, IntPredicate pred) {  // 000[1]11
  Int mid, _hi = hi;
  while (lo < hi) {
    mid = lo + (hi - lo)/2;
    if (pred(mid)) {
      hi = mid;
    } else {
      lo = mid + 1;
    }
  }
  if (!pred(lo)) {
    return _hi;  // All false.
  }
  return lo;
}

template<class Int, class IntPredicate>
Int binary_search_last_true(Int lo, Int hi, IntPredicate pred) {  // 11[1]000
  Int mid, _hi = hi--;
  while (lo < hi) {
    mid = lo + (hi - lo + 1)/2;
    if (pred(mid)) {
      lo = mid;
    } else {
      hi = mid - 1;
    }
  }
  if (!pred(lo)) {
    return _hi;  // All false.
  }
  return lo;
}

template<class DoublePredicate>
double fbinary_search(double lo, double hi, DoublePredicate pred) {  // 000[1]11
  double mid;
  for (int i = 0; i < 100; i++) {
    mid = (lo + hi)/2.0;
    if (pred(mid)) {
      hi = mid;
    } else {
      lo = mid;
    }
  }
  return lo;
}
\end{lstlisting}
\textbf{Example Usage}\par
\begin{lstlisting}[language=C++]
#include <cassert>
#include <cmath>

// Simple predicate examples.
bool pred1(int x) { return x >= 3; }
bool pred2(int x) { return false; }
bool pred3(int x) { return x <= 5; }
bool pred4(int x) { return true; }
bool pred5(double x) { return x >= 1.2345; }

int main() {
  assert(binary_search_first_true(0, 7, pred1) == 3);
  assert(binary_search_first_true(0, 7, pred2) == 7);
  assert(binary_search_last_true(0, 7, pred3)  == 5);
  assert(binary_search_last_true(0, 7, pred4)  == 6);
  assert(fabs(fbinary_search(-10.0, 10.0, pred5) - 1.2345) < 1e-15);
  return 0;
}
\end{lstlisting}

\section{Unimodal and Continuous Search}
\setcounter{section}{2}
\setcounter{subsection}{0}
\subsection{Ternary Search}
Given a unimodal function \inlinecode{f(x)} taking a single \inlinecode{double} argument, find its global maximum or minimum point to a specified absolute error.

\inlinecode{ternary\_search\_min()} takes the domain \inlinecode{[lo, hi]} of a continuous function \inlinecode{f(x)} and returns a number \inlinecode{x} such that \inlinecode{f} is strictly decreasing on the interval \inlinecode{[lo, x]} and strictly increasing on the interval \inlinecode{[x, hi]}. For the function to be correct and deterministic, such an \inlinecode{x} must exist and be unique.

\inlinecode{ternary\_search\_max()} takes the domain \inlinecode{[lo, hi]} of a continuous function \inlinecode{f(x)} and returns a number \inlinecode{x} such that \inlinecode{f} is strictly increasing on the interval \inlinecode{[lo, x]} and strictly decreasing on the interval \inlinecode{[x, hi]}. For the function to be correct and deterministic, such an \inlinecode{x} must exist and be unique.

\textbf{Time Complexity}: $O(\log n)$ calls will be made to \inlinecode{f()}, where $n$ is the distance between \inlinecode{lo} and \inlinecode{hi} divided by the specified absolute error (epsilon).

\textbf{Space Complexity}: $O(1)$ auxiliary for both operations.

\begin{lstlisting}[language=C++]
template<class UnimodalFunction>
double ternary_search_min(double lo, double hi, UnimodalFunction f,
                          const double EPS = 1e-12) {
  double lthird, hthird;
  while (hi - lo > EPS) {
    lthird = lo + (hi - lo)/3;
    hthird = hi - (hi - lo)/3;
    if (f(lthird) < f(hthird)) {
      hi = hthird;
    } else {
      lo = lthird;
    }
  }
  return lo;
}

template<class UnimodalFunction>
double ternary_search_max(double lo, double hi, UnimodalFunction f,
                          const double EPS = 1e-12) {
  double lthird, hthird;
  while (hi - lo > EPS) {
    lthird = lo + (hi - lo)/3;
    hthird = hi - (hi - lo)/3;
    if (f(lthird) < f(hthird)) {
      lo = lthird;
    } else {
      hi = hthird;
    }
  }
  return hi;
}
\end{lstlisting}
\textbf{Example Usage}\par
\begin{lstlisting}[language=C++]
#include <cassert>
#include <cmath>

bool equal(double a, double b) {
  return fabs(a - b) < 1e-7;
}

// Parabola opening up with vertex at (-2, -24).
double f1(double x) {
  return 3*x*x + 12*x - 12;
}

// Parabola opening down with vertex at (2/19, 8366/95) ~ (0.105, 88.063).
double f2(double x) {
  return -5.7*x*x + 1.2*x + 88;
}

// Absolute value function shifted to the right by 30 units.
double f3(double x) {
  return fabs(x - 30);
}

int main() {
  assert(equal(ternary_search_min(-1000, 1000, f1), -2));
  assert(equal(ternary_search_max(-1000, 1000, f2), 2.0/19));
  assert(equal(ternary_search_min(-1000, 1000, f3), 30));
  return 0;
}
\end{lstlisting}
\subsection{Hill Climbing}
Given a continuous function \inlinecode{f(x, y)} to \inlinecode{double} and a (possibly arbitrary) initial guess (\inlinecode{x0}, \inlinecode{y0}), return a potential global minimum found through hill-climbing. Optionally, two \inlinecode{double} pointers \inlinecode{critical\_x} and \inlinecode{critical\_y} may be passed to store the input points to \inlinecode{f} at which the returned minimum value is attained.

Hill-climbing is a heuristic which starts at the guess, then considers taking a single step in each of a fixed number of directions. The direction with the best (in this case, minimum) value is chosen, and further steps are taken in it until the answer no longer improves. When this happens, the step size is reduced and the same process repeats until a desired absolute error is reached. The technique's success heavily depends on the behavior of \inlinecode{f} and the initial guess. Therefore, the result is not guaranteed to be the global minimum.

\textbf{Time Complexity}: $O(d \log n)$ call will be made to \inlinecode{f}, where $d$ is the number of directions considered at each position and $n$ is the search space that is approximately proportional to the maximum possible step size divided by the minimum possible step size.

\textbf{Space Complexity}: $O(1)$ auxiliary.

\begin{lstlisting}[language=C++]
#include <cstddef>
#include <cmath>

template<class ContinuousFunction>
double find_min(ContinuousFunction f, double x0, double y0,
                double *critical_x = NULL, double *critical_y = NULL,
                const double STEP_MIN = 1e-9, const double STEP_MAX = 1e6,
                const int NUM_DIRECTIONS = 6) {
  static const double PI = acos(-1.0);
  double x = x0, y = y0, res = f(x0, y0);
  for (double step = STEP_MAX; step > STEP_MIN; ) {
    double best = res, best_x = x, best_y = y;
    bool found = false;
    for (int i = 0; i < NUM_DIRECTIONS; i++) {
      double a = 2.0*PI*i / NUM_DIRECTIONS;
      double x2 = x + step*cos(a), y2 = y + step*sin(a);
      double value = f(x2, y2);
      if (best > value) {
        best_x = x2;
        best_y = y2;
        best = value;
        found = true;
      }
    }
    if (!found) {
      step /= 2.0;
    } else {
      x = best_x;
      y = best_y;
      res = best;
    }
  }
  if (critical_x != NULL && critical_y != NULL) {
    *critical_x = x;
    *critical_y = y;
  }
  return res;
}
\end{lstlisting}
\textbf{Example Usage}\par
\begin{lstlisting}[language=C++]
#include <cassert>
#include <cmath>

bool eq(double a, double b) {
  return fabs(a - b) < 1e-8;
}

// Paraboloid with global minimum at f(2, 3) = 0.
double f(double x, double y) {
  return (x - 2)*(x - 2) + (y - 3)*(y - 3);
}

int main() {
  double x, y;
  assert(eq(find_min(f, 0, 0, &x, &y), 0));
  assert(eq(x, 2) && eq(y, 3));
  return 0;
}
\end{lstlisting}

\section{Dynamic Programming Optimization}
\setcounter{section}{3}
\setcounter{subsection}{0}
\subsection{Convex Hull Trick (Semi-Dynamic)}
Given a set of pairs $(m, b)$ specifying lines of the form $y = mx + b$, process a set of x-coordinate queries each asking to find the minimum y-value when any of the given lines are evaluated at the specified x. For each \inlinecode{add\_line(m, b)} call, \inlinecode{m} must be the minimum \inlinecode{m} of all lines added so far. For each \inlinecode{query(x)} call, \inlinecode{x} must be the maximum \inlinecode{x} of all queries made so far.

The following implementation is a concise, semi-dynamic version of the convex hull optimization technique. It supports an an interlaced sequence of \inlinecode{add\_line()} and \inlinecode{query()} calls, as long as the preconditions of descending \inlinecode{m} and ascending \inlinecode{x} are satisfied. As a result, it may be necessary to sort the lines and queries before calling the functions. In that case, the overall time complexity will be dominated by the sorting step.

\textbf{Time Complexity}: $O(n)$ for any interlaced sequence of \inlinecode{add\_line()} and \inlinecode{query()} calls, where $n$ is the number of lines added. This is because the overall number of steps taken by \inlinecode{add\_line()} and \inlinecode{query()} are respectively bounded by the number of lines. Thus a single call to either \inlinecode{add\_line()} or \inlinecode{query()} will have an amortized $O(1)$ running time.

\textbf{Space Complexity}:
\begin{compactitem}
  \item $O(n)$ for storage of the lines.
  \item $O(1)$ auxiliary for \inlinecode{add\_line()} and \inlinecode{query()}.
\end{compactitem}

\begin{lstlisting}[language=C++]
#include <vector>

std::vector<long long> M, B;
int ptr = 0;

void add_line(long long m, long long b) {
  int len = M.size();
  while (len > 1 && (B[len - 2] - B[len - 1])*(m - M[len - 1]) >=
                    (B[len - 1] - b)*(M[len - 1] - M[len - 2])) {
    len--;
  }
  M.resize(len);
  B.resize(len);
  M.push_back(m);
  B.push_back(b);
}

long long query(long long x) {
  if (ptr >= (int)M.size()) {
    ptr = (int)M.size() - 1;
  }
  while (ptr + 1 < (int)M.size() &&
         M[ptr + 1]*x + B[ptr + 1] <= M[ptr]*x + B[ptr]) {
    ptr++;
  }
  return M[ptr]*x + B[ptr];
}
\end{lstlisting}
\textbf{Example Usage}\par
\begin{lstlisting}[language=C++]
#include <cassert>

int main() {
  add_line(3, 0);
  add_line(2, 1);
  add_line(1, 2);
  add_line(0, 6);
  assert(query(0) == 0);
  assert(query(1) == 3);
  assert(query(2) == 4);
  assert(query(3) == 5);
  return 0;
}
\end{lstlisting}
\subsection{Convex Hull Trick (Fully-Dynamic)}
Given a set of pairs $(m, b)$ specifying lines of the form $y = mx + b$, process a set of x-coordinate queries each asking to find the minimum y-value when any of the given lines are evaluated at the specified x. To instead have the queries optimize for maximum y-value, call the constructor with \inlinecode{query\_max=true}.

The following implementation is a fully dynamic variant of the convex hull optimization technique, using a self-balancing binary search tree (\inlinecode{std::set}) to support the ability to call \inlinecode{add\_line()} and \inlinecode{query()} in any desired order.

\textbf{Time Complexity}: $O(\log n)$ amortized per call to \inlinecode{add\_line()} and $O(\log n)$ per call to \inlinecode{query()}, where $n$ is the number of lines added.

\textbf{Space Complexity}:
\begin{compactitem}
  \item $O(n)$ for storage of the lines.
  \item $O(1)$ auxiliary for \inlinecode{add\_line()} and \inlinecode{query()}.
\end{compactitem}

\begin{lstlisting}[language=C++]
#include <cassert>
#include <climits>
#include <set>

class hull_optimizer {
  struct line {
    long long m, b;
    mutable long long xhi;
    bool is_query;

    line(long long m, long long b, long long xhi = 0, bool is_query = false)
        : m(m), b(b), xhi(xhi), is_query(is_query) {}

    bool operator<(const line &l) const {
      return l.is_query ? xhi < l.xhi : m < l.m;
    }
  };

  std::multiset<line> hull;
  bool query_max;

  typedef std::multiset<line>::iterator hulliter;

  static long long div_floor(long long a, long long b) {
    return a/b - ((a ^ b) < 0 && a % b);
  }

  bool update_border(hulliter x, hulliter y) {
    if (y == hull.end()) {
      x->xhi = LLONG_MAX;
      return false;
    }
    if (x->m == y->m) {
      x->xhi = (x->b > y->b) ? LLONG_MAX : LLONG_MIN;
    } else {
      x->xhi = div_floor(y->b - x->b, x->m - y->m);
    }
    return x->xhi >= y->xhi;
  }

 public:
  hull_optimizer(bool query_max = false) : query_max(query_max) {}

  void add_line(long long m, long long b) {
    if (!query_max) {
      m = -m;
      b = -b;
    }
    hulliter z = hull.insert(line(m, b));
    hulliter y = z++;
    hulliter x = y;
    while (update_border(y, z)) {
      z = hull.erase(z);
    }
    if (x != hull.begin() && update_border(--x, y)) {
      update_border(x, y = hull.erase(y));
    }
    while ((y = x) != hull.begin() && (--x)->xhi >= y->xhi) {
      update_border(x, hull.erase(y));
    }
  }

  long long query(long long x) const {
    assert(!hull.empty());
    line q(0, 0, x, true);
    hulliter it = hull.lower_bound(q);
    long long res = it->m*x + it->b;
    return query_max ? res : -res;
  }
};
\end{lstlisting}
\textbf{Example Usage}\par
\begin{lstlisting}[language=C++]
#include <cassert>

int main() {
  hull_optimizer h;
  h.add_line(3, 0);
  h.add_line(0, 6);
  h.add_line(1, 2);
  h.add_line(2, 1);
  assert(h.query(0) == 0);
  assert(h.query(2) == 4);
  assert(h.query(1) == 3);
  assert(h.query(3) == 5);
  return 0;
}
\end{lstlisting}

\section{Linear and Convex Optimization}
\setcounter{section}{4}
\setcounter{subsection}{0}
\subsection{Linear Programming (Simplex)}
Solves a linear programming problem using Dantzig's simplex algorithm. The canonical form of a linear programming problem is to maximize (or minimize) the dot product $cx$, subject to $ax \leq b$ and $x \geq 0$, where $x$ is a vector of unknowns to be solved, $c$ is a vector of coefficients, $a$ is a matrix of linear equation coefficients, and $b$ is a vector of boundary coefficients.

\begin{compactitem}
  \item \inlinecode{simplex\_solve(a, b, c, \&x)} solves the linear programming problem for an $m$ by $n$ matrix \inlinecode{a} of real values, a length $m$ vector \inlinecode{b}, and a length $n$ vector \inlinecode{c}, returning 0 if a solution was found or $-1$ if there are no solutions. If a solution is found, then the vector pointed to by \inlinecode{x} is populated with the solution vector of length $n$.
\end{compactitem}

\textbf{Time Complexity}: Polynomial (average) on the number of equations and unknowns, but exponential in the worst case.

\textbf{Space Complexity}: $O(m \cdot n)$ auxiliary heap space.

\begin{lstlisting}[language=C++]
#include <cmath>
#include <limits>
#include <vector>

template<class Matrix>
int simplex_solve(const Matrix &a, const std::vector<double> &b,
                  const std::vector<double> &c, std::vector<double> *x,
                  const bool MAXIMIZE = true, const double EPS = 1e-10) {
  int m = a.size(), n = c.size();
  Matrix t(m + 2, std::vector<double>(n + 2));
  t[1][1] = 0;
  for (int j = 1; j <= n; j++) {
    t[1][j + 1] = MAXIMIZE ? c[j - 1] : -c[j - 1];
  }
  for (int i = 1; i <= m; i++) {
    for (int j = 1; j <= n; j++) {
      t[i + 1][j + 1] = -a[i - 1][j - 1];
    }
    t[i + 1][1] = b[i - 1];
  }
  for (int j = 1; j <= n; j++) {
    t[0][j + 1] = j;
  }
  for (int i = n + 1; i <= m + n; i++) {
    t[i - n + 1][0] = i;
  }
  double p1 = 0, p2 = 0;
  bool done = true;
  do {
    double mn = std::numeric_limits<double>::max(), xmax = 0, v;
    for (int j = 2; j <= n + 1; j++) {
      if (t[1][j] > 0 && t[1][j] > xmax) {
        p2 = j;
        xmax = t[1][j];
      }
    }
    for (int i = 2; i <= m + 1; i++) {
      v = fabs(t[i][1] / t[i][p2]);
      if (t[i][p2] < 0 && mn > v) {
        mn = v;
        p1 = i;
      }
    }
    std::swap(t[p1][0], t[0][p2]);
    for (int i = 1; i <= m + 1; i++) {
      if (i != p1) {
        for (int j = 1; j <= n + 1; j++) {
          if (j != p2) {
            t[i][j] -= t[p1][j]*t[i][p2] / t[p1][p2];
          }
        }
      }
    }
    t[p1][p2] = 1.0 / t[p1][p2];
    for (int j = 1; j <= n + 1; j++) {
      if (j != p2) {
        t[p1][j] *= fabs(t[p1][p2]);
      }
    }
    for (int i = 1; i <= m + 1; i++) {
      if (i != p1) {
        t[i][p2] *= t[p1][p2];
      }
    }
    for (int i = 2; i <= m + 1; i++) {
      if (t[i][1] < 0) {
        return -1;
      }
    }
    done = true;
    for (int j = 2; j <= n + 1; j++) {
      if (t[1][j] > 0) {
        done = false;
      }
    }
  } while (!done);
  x->clear();
  for (int j = 1; j <= n; j++) {
    for (int i = 2; i <= m + 1; i++) {
      if (fabs(t[i][0] - j) < EPS) {
        x->push_back(t[i][1]);
      }
    }
  }
  return 0;
}
\end{lstlisting}
\textbf{Example Usage}\par
\begin{lstlisting}[language=C++]
#include <cassert>
#include <iostream>
using namespace std;

int main() {
  // Solve [x, y] that maximizes 3x + 4y, subject to x, y >= 0 and:
  //  -2x +    1y <=  0
  //   1x + 0.85y <=  9
  //   1x +    2y <= 14
  const int equations = 3, unknowns = 2;
  double a[equations][unknowns] = {{-2, 1}, {1, 0.85}, {1, 2}};
  double b[equations] = {0, 9, 14};
  double c[unknowns] = {3, 4};
  vector<vector<double>> va(equations, vector<double>(unknowns));
  vector<double> vb(b, b + equations), vc(c, c + unknowns), x;
  for (int i = 0; i < equations; i++) {
    for (int j = 0; j < unknowns; j++) {
      va[i][j] = a[i][j];
    }
  }
  assert(simplex_solve(va, vb, vc, &x) == 0);
  double maxval = 0;
  for (int i = 0; i < (int)x.size(); i++) {
    maxval += c[i]*x[i];
  }
  cout << "Solution = " << maxval << " at (" << x[0];
  for (int i = 1; i < (int)x.size(); i++) {
    cout << ", " << x[i];
  }
  cout << ")." << endl;
  return 0;
}
\end{lstlisting}
\begin{exampleoutput}
Solution = 33.3043 at (5.30435, 4.34783).
\end{exampleoutput}
